\clearpage
\section{よくある質問・困ったとき}\label{sec:faq}

\newcommand\faq[2]{\paragraph{Q.\ #1}\mbox{}\\A.\ #2\par\medskip}

\faq{アプリが開けません（「開いていません」「開けません」と出る）。}{
MacStarStacker は Apple の Developer ID による署名・公証をしていないため、初めて開くときに止められます。
DMG の「かんたんインストーラ」を使うか、\ref{sec:install-manual}節の手順で開いてください。}

\faq{プレビューが真っ暗で、何が写っているか分かりません。}{
プレビューの上の\ui{Auto Stretch}にチェックを入れると、見やすい明るさで表示します（表示だけで、結果は変わりません）。}

\faq{合成した結果が暗い（または色が変）です。}{
RAW (DNG) で書き出した場合は、現像前のデータなので、Lightroom などで露出・ホワイトバランスを調整してください。元の RAW と同じように調整できます。}

\faq{「…のメモリが必要なため実行できません」と出ました。}{
中央値合成・シグマクリッピング・新星景モードは、全部の写真を一度にメモリに置きます。写真の枚数を減らすか、中央値の代わりに平均を使ってください。}

\faq{新星景モードで「空と地上を分けずに合成しました」と出ました。}{
メッセージに理由が表示されます。星がほとんど動いていない（撮影時間が短い）、空だけ・地上だけの構図、などです。\ref{sec:nightscape}節の「うまくいかないとき」を見てください。}

\faq{星の位置合わせに失敗しました、と出ます。}{
雲が多い、星が少ない（明るすぎる・ピントが合っていない）写真があると、位置を合わせられないことがあります。その写真をファイルリストから外して（\ui{×}）やり直してください。}

\faq{JPEG の写真でも使えますか？}{
使えます。ただし RAW で撮った写真の方が、合成した後に明るさや色を大きく調整しても画質が落ちにくいのでおすすめです。}

\faq{古い Mac（Metal に対応していない Mac）でも使えますか？}{
\textsf{macOS10.13} 版では、Metal（GPU）が使えない Mac では CPU で合成します。時間はかかりますが、結果は同じです。
また、RAW のプレビューは読み込みに時間がかかることがあります。}

\faq{作業を最初からやり直したいです。}{
ファイルリストの右上の\ui{すべてクリア}（または\menu{編集 → すべてクリア…}）で、起動した直後の状態に戻せます。}
